\documentclass{article}

\usepackage{amsmath,amssymb}
\usepackage{xurl}
\usepackage[hidelinks]{hyperref}
\setlength{\emergencystretch}{2em}

% Shared commands for notes in docs/paper-gaps/.

\DeclareUnicodeCharacter{2080}{\ifmmode{}_{0}\else\textsubscript{0}\fi}
\DeclareUnicodeCharacter{2081}{\ifmmode{}_{1}\else\textsubscript{1}\fi}
\DeclareUnicodeCharacter{2082}{\ifmmode{}_{2}\else\textsubscript{2}\fi}
\DeclareUnicodeCharacter{2099}{\ifmmode{}_{n}\else\textsubscript{n}\fi}

\newcommand{\Lean}{\textsc{Lean}}
\ExplSyntaxOn
\NewDocumentCommand{\leanid}{m}{
  \ifmmode
    \textsf{\detokenize{#1}}
  \else
    \group_begin:
    \urlstyle{sf}
    \tl_set:Nn \l_tmpa_tl {#1}
    \tl_replace_all:Nnn \l_tmpa_tl {\_} {_}
    \exp_args:NV \path \l_tmpa_tl
    \group_end:
  \fi
}
\ExplSyntaxOff
\providecommand{\tr}{\operatorname{tr}}
\newcommand{\tnleanIssueBase}{https://github.com/LionSR/TNLean/issues/}
\newcommand{\tnleanPrBase}{https://github.com/LionSR/TNLean/pull/}
\newcommand{\tnleanDocBase}{https://sirui-lu.com/TNLean/docs/}
\newcommand{\ghissue}[1]{\href{\tnleanIssueBase#1}{GitHub issue \##1}}
\newcommand{\ghpr}[1]{\href{\tnleanPrBase#1}{PR \##1}}
\newcommand{\doclink}[1]{\href{\tnleanDocBase#1.html}{\path{#1}}}

% Dirac notation and the identity operator, as in blueprint/src/macros/common.tex.
% \providecommand keeps note-local definitions authoritative.
\usepackage{dsfont}
\providecommand{\ket}[1]{|#1\rangle}
\providecommand{\bra}[1]{\langle #1|}
\providecommand{\braket}[2]{\langle #1 | #2 \rangle}
\providecommand{\ketbra}[2]{|#1\rangle\!\langle #2|}
\providecommand{\Id}{\mathds{1}}
\providecommand{\one}{\mathds{1}}
\providecommand{\C}{\mathbb{C}}
\providecommand{\MN}[1]{M_{#1}(\C)}
\providecommand{\MD}{M_D(\C)}

% Verdict marker: \gapnote{<kind>}{<status>}. Kind is the policy
% classification (clarification, local-correction, scope-restriction,
% unfaithful, false-source, open-gap); status is open, wip (elimination
% underway), resolved, or historical. Machine-read by the site index;
% prints a verdict line.
\newcommand{\gapnote}[2]{\par\noindent\textbf{Verdict:} #1 (#2).\par}


% Fallbacks keep this note compilable when a different command.tex is found first.
\providecommand{\leanid}[1]{\textsf{\detokenize{#1}}}
\providecommand{\gapnote}[2]{\par\noindent\textbf{Verdict:} #1 (#2).\par}

\title{Nondegenerate Rectangles and Triangles in Area-Law Templates}
\author{}
\date{2026-10-08}

\begin{document}
\maketitle
\gapnote{local-correction}{resolved}

\section{The source definition}

Definition~9.3 of the two-dimensional area-law manuscript of September~24,
2026~\cite{OpenAI2026AreaLaw} (\path{08-scanner.tex}, lines~546--564) defines
a template with parameters $(n,s_0)$ as a nonempty finite set
$T\subset\mathbb Z^2$ of the form
\[
  T=\bigcup_{a=1}^{N_0}(P_a\cap\mathbb Z^2),
  \qquad n\ge C_{\rm tpl}N_0(s_0+1),
  \qquad \operatorname{diam}_\infty T\le n,
\]
with $N_0\ge1$, where each $P_a$ is a closed convex rectangle or triangle of
sup-norm diameter at most $s_0$, all of whose sides have slope
$0,\infty,1$ or $-1$.

\section{The convention}

The words rectangle and triangle are read in their ordinary sense: a triangle
has three noncollinear vertices, and a rectangle has two nonzero orthogonal
adjacent sides. A literal reading that also admitted a segment or a point as a
degenerate rectangle or triangle is not intended by the source. Its remark
after the definition (lines~566--569) states that the estimates use the
pieces only through their number $N_0$ and their sizes. The nonemptiness of
$T$ and the bound $N_0\ge1$ are printed in the source and are not part of the
convention.

The convention is built into the definition of the permitted
polygons.\footnote{\leanid{TNLean.PEPS.AreaLaw.Geometry.TemplatePolygon}.}
No downstream statement carries a separate nondegeneracy hypothesis.

\section{Effect}

The convention is a reading of the source's words, not an additional
hypothesis on any theorem, and no elimination is planned.

\bibliographystyle{alpha}
\bibliography{references}

\end{document}
